%=====================================================================
\section{Violated assumptions, duplicated definitions, absent terms}
\label{sec:defects}

This section is part of the model specification, not an appendix to it. Every
item below is either an assumption the operating regime violates, a physical
quantity the code defines twice, or a term omitted altogether. A reader
reproducing these results needs all three.

%---------------------------------------------------------------------
\subsection{H3 is saturated, and it now carries a design decision}
\label{sec:defects:h3}

\subsubsection*{The limit was being tested at the wrong radius}

Hypothesis~H3 takes the melt region to be a concentric annulus whose front never
meets the front of a neighbouring leg. Until this version the geometric check
guarding it compared $r_{e}+\delta$ against the \emph{borehole wall}. That is
the wrong radius. The $2n_{t}=4$ legs share the borehole cross-section equally,
so each owns
\begin{equation}
  \pi r_{\mathrm{cell}}^{2} = \frac{\pi R^{2}}{2n_{t}},
  \qquad\text{i.e.}\qquad
  r_{\mathrm{cell}} = \frac{R}{\sqrt{2n_{t}}},
  \qquad
  \delta_{\mathrm{merge}} = r_{\mathrm{cell}} - r_{e},
  \label{eq:rcell}
\end{equation}
with $R=D_{\mathrm{well}}/2$. Numerically $r_{\mathrm{cell}}=\SI{44.45}{\milli\meter}$
against a wall at \SI{88.90}{\milli\meter}, so the check permitted
$\delta=\SI{67.82}{\milli\meter}$ where the true limit is
\SI{23.37}{\milli\meter} --- a factor \num{2.9} too permissive, and therefore
silent at every design point ever run. Equation~\eqref{eq:rcell} is now
\code{Case.r\_cell} and \code{Case.delta\_merge}, and the plots, the diagnostics
and the sizing report all read the same limit.

\subsubsection*{Under Formulation C the limit cannot be exceeded --- it is sat on}

Correcting the radius turns out not to produce a violation, for a reason worth
stating precisely:
\begin{equation}
  A\!\left(\delta_{\mathrm{merge}}\right) \equiv A_{\mathrm{avail}},
\end{equation}
identically --- a leg owns exactly the PCM inside its own cell, so melting all
of it puts the front exactly on the cell boundary. Since Formulation~C caps
$\Amelt$ at $A_{\mathrm{avail}}$ (\S\ref{sec:front:enthalpy}), it follows that
$\delta \le \delta_{\mathrm{merge}}$ always. \emph{$\varepsilon_{\mathrm{local}}=1$
and ``the fronts merge'' are the same line.}

What the corrected diagnostic reveals instead is that the design \emph{sits} on
that line: at the v0.4 design point $\delta/\delta_{\mathrm{merge}}=1.000$ over
\SI{99}{\percent} of the well. The annular conduction resistance is being
evaluated at the exact limit of its validity almost everywhere. Two errors
follow, and both point the same way:

\begin{enumerate}[leftmargin=1.7em,itemsep=3pt]
\item The resistance $\ln\!\left(1+\delta/r_{e}\right)/2\pi\km$ assumes the full
  azimuth $2\pi\left(r_{e}+\delta\right)$ conducts. Past merge it does not --- a
  neighbour has taken part of it --- and the residual solid sits in the
  \emph{corners} between legs, on a longer and narrower path. Late-stage
  resistance is \textbf{understated}.
\item Fins enter $\Ui$ only as added \emph{surface at the tube}, through
  $\etaf$ and $P_{T}$ (\S\ref{sec:network:fins}). They are not represented as
  radial conduction paths
  reaching into that corner PCM --- which is the single thing a fin does best in
  this regime. The benefit is \textbf{uncredited}.
\end{enumerate}

\subsubsection*{Why this now matters: the fin count}

This ceased to be a bookkeeping concern when it began to decide geometry.
Sweeping the fins gives:

\begin{center}
\begin{tabular}{@{}l r r r l@{}}
\toprule
configuration & $N_{\mathrm{capacity}}$ & $N_{\mathrm{rate}}$
              & $\Nwells$ & binding \\
\midrule
$24\times\SI{7.50}{\milli\meter}$ & 11.573 & 11.263 & \textbf{11.573} & capacity \\
$16\times\SI{7.50}{\milli\meter}$ & 11.348 & 11.287 & \textbf{11.348} & capacity \\
$12\times\SI{7.50}{\milli\meter}$ & 11.239 & 11.529 & \textbf{11.529} & rate \\
$\phantom{0}8\times\SI{7.50}{\milli\meter}$ & 11.132 & 12.208 & \textbf{12.208} & rate \\
bare tube & 10.924 & 16.569 & \textbf{16.569} & rate \\
\bottomrule
\end{tabular}
\end{center}

\noindent
Read the $N_{\mathrm{capacity}}$ column alone and the conclusion is that fins
make matters worse --- it falls monotonically as fin metal is removed. That
column contains \emph{no heat transfer at all}: it is
$\Delta E/\left(\rho V_{\mathrm{well}}\left[\hm+c_{p,l}\Delta T\right]\right)$,
and removing a fin returns its metal volume to the PCM. Indeed
$N_{\mathrm{capacity}}\cdot V_{\mathrm{well}}$ is constant to machine precision
across the whole sweep. The real answer, $\Nwells=\max$ of the two criteria,
moves the other way: removing the fins costs \SI{42}{\percent} more wells.
\S\ref{sec:hydraulics:capacity} states the reporting rule that follows.

The optimum sits where the two criteria cross, at $16\times\SI{7.5}{\milli\meter}$
--- so the design point is mildly over-finned, by about \SI{2}{\percent} in well
count. But given the two biases above, that figure should be read as a
\textbf{floor} on the fin count rather than an optimum. Settling it requires a
two-dimensional conduction calculation in the cell, which is the most
consequential piece of open work in this model.

\subsubsection*{The charging window is what decides it}

Fins buy \emph{rate}, so their worth is set by how much rate is demanded:

\begin{center}
\begin{tabular}{@{}r r r r@{}}
\toprule
$t_{ch}$ [h] & bare & $24$ fins & gain \\
\midrule
 4 & 28.976 & 15.794 & 1.84 \\
 6 & 22.579 & 12.499 & 1.81 \\
10 & 16.569 & 11.573 & 1.43 \\
14 & 13.468 & 11.573 & 1.16 \\
20 & 11.432 & 11.573 & 0.99 \\
\bottomrule
\end{tabular}
\end{center}

\noindent
Break-even is near \SI{20}{\hour}. The origin is visible in the closed form of
\S\ref{sec:front:legacy}: putting $x=r_{\mathrm{cell}}/r_{e}$ in
Eq.~\eqref{eq:stefan-implicit}, a \emph{bare} tube at $\Delta T=\SI{10}{\kelvin}$ reaches
the cell boundary in
\begin{equation}
  t = \frac{\rhom\hm r_{e}^{2}}{\km\Delta T}
      \left[\tfrac12 x^{2}\ln x - \tfrac14 x^{2} + \tfrac14\right]
    \approx \SI{6.4}{\hour},
\end{equation}
comfortably inside the \SI{10}{\hour} window --- which is why a bare tube
\emph{nearly} works here, and why the fins look marginal. Halve the window and
they are not marginal at all.

%---------------------------------------------------------------------
\subsection{H9 --- azimuthal insulation between cascade layers}
\label{sec:defects:h9}

\emph{Assumed perfect in v0.4a so that the rest of the model can be evaluated.
This section states what is being assumed away, and what it is worth, so the
assumption can be tested rather than inherited.}

\subsubsection*{Where the requirement comes from}

The cascade (H8) assigns melting temperatures along the \emph{developed} tube
length $L_{\mathrm{tube}}=2L_{\mathrm{well}}$, because that is the coordinate
the fluid actually traverses. A hairpin's down leg occupies
$s\in[0,L_{\mathrm{well}}]$ and its return leg $s\in[L_{\mathrm{well}},
L_{\mathrm{tube}}]$, so by Eq.~\eqref{eq:depth-fold} the two legs at a given
depth $d$ sit at developed positions $s=d$ and $s=L_{\mathrm{tube}}-d$ --- in
different layers:

\begin{center}
\begin{tabular}{@{}r r r r@{}}
\toprule
depth [m] & down leg $\Tm$ [\degC] & return leg $\Tm$ [\degC] & gap [K] \\
\midrule
   0 & 150.0 & 101.1 & \textbf{48.9} \\
 400 & 143.9 & 107.2 & 36.7 \\
 800 & 137.8 & 113.3 & 24.4 \\
1200 & 131.7 & 119.4 & 12.2 \\
1524 & 125.6 & 125.6 & \phantom{0}0.0 \\
\bottomrule
\end{tabular}
\end{center}

\noindent
The borehole cross-section must therefore be \emph{partitioned}: at every depth
the PCM around the down legs is a different material from the PCM around the
return legs, and the two are asked to hold a temperature difference of up to
\SI{48.9}{\kelvin} across a gap of a few centimetres. The design is silent on
how this partition is built.

The per-leg area bookkeeping is nonetheless correct.
$A_{\mathrm{avail}}=V_{\mathrm{well}}/(n_{t}L_{\mathrm{tube}})$ and
$L_{\mathrm{tube}}=2L_{\mathrm{well}}$, so the divisor is the total \emph{leg}
length $2n_{t}L_{\mathrm{well}}$ and each of the $2n_{t}=4$ legs owns exactly a
quarter of the PCM cross-section, \SI{4.5411e-3}{\meter\squared}. That is what
Eq.~\eqref{eq:rcell} encodes, and it is unaffected by how the partition is
drawn --- a quarter-disc and a half-of-a-half-disc have the same area, hence the
same $\delta_{\mathrm{merge}}$.

\subsubsection*{What it is worth}

Treating the interface between adjacent cells as a slab of width and path length
both $\sim 2r_{\mathrm{cell}}$ gives $q'_{\mathrm{leak}} \sim \km\,\Delta T$:

\begin{center}
\begin{tabular}{@{}l r r@{}}
\toprule
 & $\Delta T$ [K] & $q'_{\mathrm{leak}}$ [\si{\watt\per\meter}] \\
\midrule
wellhead    & 48.9 & 22.0 \\
mid-depth   & 24.4 & 11.0 \\
well bottom &  0.0 &  0.0 \\
\midrule
useful duty $q'$, time-mean over the charge & & 82.8 \\
\bottomrule
\end{tabular}
\end{center}

\noindent
Of order \SI{25}{\percent} of the useful duty at the wellhead. Even allowing a
shape factor two or three times kinder than a slab, this is a first-order term,
and its sign is unfavourable in a specific way: the leak flows from the hot
layer to the cold layer, which is exactly the temperature separation the cascade
is built to create. An uninsulated partition does not merely lose energy --- it
partially undoes $N_{\mathrm{lay}}$.

Three consequences follow for the work still to be done.

\begin{enumerate}[leftmargin=1.7em,itemsep=3pt]
\item \textbf{$N_{\mathrm{lay}}$ can no longer be swept for free.} More layers
  means a finer cascade and a better match to the glide, but also a larger
  layer-to-layer $\Delta T$ at the wellhead and hence a larger leak. There is an
  optimum, and the present model cannot see it because the leak is zero by
  assumption.
\item \textbf{The neighbour relationship is not uniform.} Within one half of the
  partition a leg's neighbour carries the \emph{same} $\Tm$, so their fronts
  meeting is a genuine merge (H3, \S\ref{sec:defects:h3}). Across the partition
  the neighbour carries a different $\Tm$, so the same geometric event is a heat
  leak, not a merge. The model treats all four legs identically and adiabatically.
\item \textbf{The partition displaces PCM.} Whatever insulates the two halves
  occupies borehole volume and therefore reduces $V_{\mathrm{well}}$, which by
  \S\ref{sec:hydraulics:capacity} raises $N_{\mathrm{capacity}}$ directly. The
  insulation is not free even before its thermal performance is considered.
\end{enumerate}

\subsubsection*{How to settle it}

The same two-dimensional conduction calculation in the cell that
\S\ref{sec:defects:h3} requires would also deliver this, since both are
questions about what happens at a cell's outer boundary. One calculation, on the
real cross-section rather than an equivalent annulus, resolves the corner PCM,
the reduced conducting azimuth, the fins' reach into the corners, and the
cross-partition leak together. Until then H9 is held as an idealisation and the
well counts in this document should be read as optimistic.

%---------------------------------------------------------------------
\subsection{Where the closed-form front fails}

Formulation~A (\S\ref{sec:front:legacy}) rests on two premises, both false in
this application, acting in opposite directions.

\paragraph{Surface mismatch.} The fluid gives up heat through the finned
perimeter $P_{T}=\SI{0.4924}{\meter}$ of Eq.~\eqref{eq:perim}; the front absorbs
it through the bare perimeter $2\pi r_{e}=\SI{0.1324}{\meter}$ assumed by
Eq.~\eqref{eq:stefan-explicit}. Factor $3.72$. The front therefore
\emph{under}-melts.

\paragraph{Retroactive driving temperature.}
Equation~\eqref{eq:stefan-explicit} assumes $\Tre$ constant since $t=0$, and is
handed its current value. Measured at mid-well over the charging window,
$\Tre-\Tm$ runs $0 \to \SI{8.62}{\kelvin}$. Applying the final value
retroactively \emph{over}-melts by a factor of about $1.56$.

The two are separable by setting $n_{f}=0$, which makes the surfaces coincide
and leaves only the second error:

\begin{center}
\begin{tabular}{@{}c c c@{}}
\toprule
$n_{f}$ & $k_{w}$ [\si{\watt\per\meter\per\kelvin}]
        & $\int q'\mathrm{d}t \,/\, \rhom\hm\Vmelt$ \\
\midrule
24 & 0.45 & 1.032 \\
24 & 45   & 2.021 \\
 0 & 0.45 & 0.642 \\
 0 & 45   & 0.642 \\
\bottomrule
\end{tabular}
\end{center}

With the fins removed the imbalance does not vanish --- it inverts, and becomes
almost independent of $k_{w}$. The near-unity entry in the first row is
therefore not evidence of a sound model. It is a \emph{triple} cancellation
between the surface mismatch, the retroactive $\Tre$, and the conductivity error
below.

\emph{Evidence:} \code{verification/verify\_fins\_off.py},
\code{verification/verify\_frozen\_Tre.py}.

%---------------------------------------------------------------------
\subsection{The wall conductivity argument --- corrected in v0.3}
\label{sec:defects:kw}

\noindent\fbox{\parbox{0.97\linewidth}{\textbf{Correction notice.} Versions of
this model up to and including the IHTC paper used the wrong value of the wall
thermal conductivity during charging. The results reported in this document for
v0.3 and later use the corrected value.}}

\medskip
$k_{w}$ carries two physical meanings, tube-wall conductivity in
Eq.~\eqref{eq:R-wall} and fin conductivity in Eq.~\eqref{eq:fin-eff}, through a
single positional argument. In the charging chain that slot received the PCM
liquid conductivity, $\SI{0.45}{\watt\per\meter\per\kelvin}$, instead of steel
at $\SI{45}{\watt\per\meter\per\kelvin}$ --- two orders of magnitude low. The
error is visible in the original solver signature, whose parameter is named
\code{k\_m\_l}.

The consequence was not confined to the wall resistance, which is small either
way. Through Eq.~\eqref{eq:fin-eff} it collapsed the fin efficiency to
$\etaf\approx 0.33$, throttling the fluid-side conductance to roughly what a
bare cylinder can absorb, and so accidentally compensating the surface mismatch
above. This is why correcting $k_{w}$ \emph{alone} makes the model worse, and
why it could only be corrected together with the front formulation.

From v0.3 the correct value is the default
(\code{Case.charge\_uses\_wall\_conductivity = True}). Setting it \code{False}
reproduces the published v0.1 configuration, which is how the verification of
\S\ref{sec:defects:verification} is performed. Both switches must be set
together: the closed-form front \emph{and} the conductivity error. Setting only
one gives a model that never existed.

%---------------------------------------------------------------------
\subsection{Quantities defined twice}

\begin{description}[leftmargin=1.6em,itemsep=3pt]
\item[Melt-layer conductance $h_{e}$.] Equation~\eqref{eq:he} is the cylindrical
  shell used in \code{compute\_T\_re}. The small-$\delta$ branch of
  \code{compute\_U\_i} instead falls back to the plane-slab form
  $h_{e}=\km/\delta$. Same physical quantity, two formulas, iterated against
  each other in the fixed-point loop of \S\ref{sec:network:Tre}. The two agree
  as $\delta/r_{e}\to0$ and diverge as the layer grows.
\item[$E_{\mathrm{well}}$.] A \emph{capacity} per well in
  Eq.~\eqref{eq:E-well}, a \emph{duty} per well in \S\ref{sec:hydraulics:kpi}.
  Different keys in the code, one name, and related by the model input
  $1/(1+\lambda)$ --- Eq.~\eqref{eq:Ewell-ratio}.
\item[$k_{w}$.] Tube-wall conductivity in Eq.~\eqref{eq:R-wall}, fin
  conductivity in Eq.~\eqref{eq:fin-eff}, one positional argument
  (\S\ref{sec:defects:kw}). This one has already caused a
  two-order-of-magnitude error that survived to publication.
\item[Laminar Nusselt number.] $4.36$ (uniform flux) in \code{well\_marched},
  $3.66$ (uniform temperature) in the legacy march. For a melting boundary the
  isothermal value is the defensible one.
\item[Charging pressure-drop end states.] The charging drop is evaluated between
  $T_{4c}$ and $T_{4a}$, a source-loop pair, whereas $T_{2c}$ is the borehole
  return. The affected property is the water density, so the error is small, but
  the temperatures are not the ones the flow actually sees.
\end{description}

%---------------------------------------------------------------------
\subsection{Absent terms}

\begin{description}[leftmargin=1.6em,itemsep=3pt]
\item[Azimuthal conduction between cells.] Hypothesis H9,
  \S\ref{sec:defects:h9}. The largest single omission: neighbouring legs are up
  to \SI{48.9}{\kelvin} apart by construction and the term is simply absent from
  the balance of \S\ref{sec:front:segeq}.
\item[Fluid thermal inertia.] Hypothesis H7, quantified in
  \S\ref{sec:front:segeq}. Dropping $\partial T/\partial t$ leaves an
  unresolved startup transient of about \SI{38}{\minute} per half-cycle,
  \SI{6}{\percent} of the window, though only \SI{1.2}{\percent} of the latent
  inventory is involved.
\item[Formation heat loss.] Hypothesis H6. The borehole wall is adiabatic. Over
  a \SI{10}{\hour} cycle this is defensible; over seasonal storage it is not.
  The geothermal gradient is carried in the configuration and used in no energy
  balance.
\item[Natural convection in the melt.] Hypothesis H4. Conduction-only transport
  in the liquid PCM under-predicts the melting rate once the layer is
  established, and does so increasingly as $\delta$ grows. Note that this is the
  one absent term that argues \emph{against} fins rather than for them.
\item[Sensible heat in the PCM.] \emph{Supplied in v0.4} --- see
  \S\ref{sec:front:enthalpy}. What remains absent is any \emph{radial
  resolution} of temperature within a phase: each segment carries one lumped
  $T_{\mathrm{pcm}}$, so superheat appears uniformly across the liquid annulus
  rather than diffusing outward from the tube.
\item[Buoyancy head.] Hypothesis H7. The \SI{1524}{\meter} loop carries roughly
  \SI{\pm 7.5}{\bar} of static head, asymmetric between charge and discharge,
  and it is absent from Eq.~\eqref{eq:dp}.
\item[Isentropic efficiencies.] \S\ref{sec:cycles:states}. Both cycles are
  evaluated between prescribed state points, so $\eta_{\mathrm{ORC}}$ and
  $\mathrm{COP}$ are idealised.
\item[Volume change on melting.] \S\ref{sec:front:inventory}. The
  \SI{6.9}{\percent} expansion is not represented.
\item[Melt-front shape around the fins.] \S\ref{sec:front:finmetal}. The fin
  \emph{metal} is excluded from $\Amelt$ as of v0.3, but the front is still
  taken as a circular annulus; physically it bulges along the fins.
\item[Front interaction in the corners between legs.] Hypothesis H3, treated at
  length in \S\ref{sec:defects:h3}. The geometric limit is now tested at the
  cell radius rather than the borehole wall, and under Formulation~C it cannot
  be exceeded --- but it is \emph{saturated} over \SI{99}{\percent} of the well,
  so the annular resistance runs at the limit of its validity and neither the
  reduced conducting azimuth nor the fins' reach into the corner PCM is
  represented. The most consequential open item in the model.
\item[$\lambda$ as an output.] \S\ref{sec:cycles:lambda}. The storage loss is
  still an input, at \SI{5}{\percent}, while the model computes
  $\eta_{\mathrm{storage}}=0.9364$.
\end{description}

%---------------------------------------------------------------------
\subsection{Error handling}

Both v0.1 bisections return the bracket bound when they fail to converge and
report it as an answer, and the driver cell continues past it. In
\code{thums\_core} these raise \code{ThumsConvergenceError} carrying the last
iterate. Similarly, failed property lookups formerly routed into the
$\delta\to0$ branch of Eq.~\eqref{eq:he}, so that a failed solve was reported as
\emph{maximum} heat transfer; they now raise \code{ThumsPropertyError}.

Retired output names follow the same principle. \code{N\_inventory} and
\code{N\_heat} raise a \code{KeyError} carrying an explanation rather than
returning a number, and the guard overrides \code{.get()} as well as
\code{[\,]}, because a silent \code{None} from \code{.get()} is the exact
failure the guard exists to prevent (\S\ref{sec:hydraulics:twoq}).

The value $h_{e}=10^{9}$ in that branch is itself correct --- a zero-thickness
melt layer has zero conduction resistance --- notwithstanding an inverted
comment in the source and an earlier audit that repeated it.

%---------------------------------------------------------------------
\subsection{Verification status}
\label{sec:defects:verification}

\begin{description}[leftmargin=1.6em,itemsep=3pt]
\item[Regression.] \code{Case(front="closed\_form")} reproduces the frozen v0.1
  fixture to $4\times10^{-11}$ on every index; the residual is CoolProp and
  NumPy version drift.
\item[Energy closure.] Structural under Formulations~B and~C, and therefore
  \emph{not} an independent check (\S\ref{sec:front:marched}).
\item[Analytical.] Equation~\eqref{eq:stefan-explicit} verified against its own
  defining quadrature to a residual of $10^{-15}$.
\item[Numerical.] The branchwise-exact enthalpy integration of
  Eq.~\eqref{eq:enth-exact} agrees with a \num{200000}-substep explicit
  reference to $3\times10^{-13}$, including across branch crossings.
\item[Grid.] Not converged and not quoted as such: $\varepsilon_{\mathrm{PCM}}$
  moves $-0.3\,\%$ from $n_{\mathrm{times}}=40$ to $320$, one-sided, so the
  default carries a small known bias.
\item[External validation.] None. No experimental or independently computed data
  set has been identified against which any formulation can be checked. This
  is the principal gap in the verification ladder, and no amount of internal
  consistency substitutes for it.
\end{description}

%---------------------------------------------------------------------
\subsection{The ORC evaporator ran heat uphill}
\label{sec:defects:orcpinch}

Every defect above concerns the store. This one and the next concern the two
power cycles, and they are the most consequential in the document: together they
move $\eta_{\mathrm{RTE}}$ by \SI{-28}{\percent} and the well count by
\SI{+35}{\percent}.

The ORC evaporating temperature was set from the hot end of the water alone,
\begin{equation}
  T_{3e} \;=\; T_{2d} - \Delta T_{2D,3E},
  \label{eq:hot-end-rule}
\end{equation}
and nothing ever examined the rest of the exchanger. That is safe only when the
two streams have similar heat-capacity curves. Here they could hardly be less
alike. The water \emph{glides} \SI{55}{\kelvin}, from \SI{146.1}{\celsius} to
\SI{91.1}{\celsius}. The cyclopentane takes \SI{76}{\percent} of its heat
\emph{boiling at one temperature}, \SI{134.1}{\celsius}. Composite curves of
that shape cross, and these crossed comprehensively:

\begin{center}
\begin{tabular}{@{}r r r r@{}}
\toprule
$Q/Q_{\mathrm{tot}}$ & $T_{\mathrm{water}}$ [\si{\celsius}]
 & $T_{\mathrm{ORC}}$ [\si{\celsius}] & gap [\si{\kelvin}] \\
\midrule
0.000 &  91.1 & 103.9 & $-12.7$ \\
0.235 & 104.1 & 133.7 & $\mathbf{-29.5}$ \\
0.700 & 129.7 & 134.1 & $-4.4$ \\
1.000 & 146.1 & 134.1 & $+12.0$ \\
\bottomrule
\end{tabular}
\end{center}

\noindent
Only the top fifth of the surface had heat flowing in the right direction. The
reported $\eta_{\mathrm{ORC}}=0.2320$ was therefore not optimistic but
\textbf{unreachable}: no isentropic efficiency and no exchanger area buys it,
because over most of the evaporator the heat was flowing from cold to hot.

\medskip
\noindent\textbf{Why it survived.} $\Delta T_{2D,3E}=\SI{12}{\kelvin}$ is a
perfectly sensible approach \emph{at the hot end}, and the hot end is the one
place anybody checks. The violation is a full \SI{30}{\kelvin} away from the
point the parameter names.

\medskip
\noindent\textbf{The repair.} \code{orc\_pinch} constructs both composite curves
properly --- real enthalpy on both sides, the latent plateau as a vertical
segment rather than a straight line in $T$, an even grid in $Q$ that includes
every break point --- and returns the minimum gap wherever it falls.
\code{feasible\_rankine} keeps Eq.~\eqref{eq:hot-end-rule} when it already
clears $\Delta T_{\mathrm{pinch,ORC}}$, and otherwise bisects $T_{3e}$ downward
until the pinch is met. The new field is
$\Delta T_{\mathrm{pinch,ORC}}=\SI{5}{\kelvin}$.

\begin{center}
\begin{tabular}{@{}l r r@{}}
\toprule
 & v0.8 & v0.9 \\
\midrule
$T_{3e}$ at CSS [\si{\celsius}] & 136.53 & 95.89 \\
$\eta_{\mathrm{ORC}}$           & 0.2320 & 0.1738 \\
$\Nwells$                       & 11.632 & 15.659 \\
$\eta_{\mathrm{RTE}}$           & 0.4367 & 0.3140 \\
\bottomrule
\end{tabular}
\end{center}

\medskip
\noindent\textbf{The design consequence, which outlives the bug.} The glide
$\Delta T_{3C,2C}$ is not a free parameter to be raised for storage density.
Raising it buys PCM inventory per well and destroys the ORC in the same
movement, because a fluid boiling at a single temperature cannot follow a
gliding one. That trade-off was invisible while the pinch went unchecked and it
is now the binding constraint in any parametric study. Buying it back requires a
zeotropic mixture, a second pressure level, or a supercritical cycle --- not a
larger exchanger.

%---------------------------------------------------------------------
\subsection{The heat pump created energy in IHX-2}
\label{sec:defects:ihx2}

The second internal heat exchanger of the high-temperature heat pump was closed
with an effectiveness expression carrying the wrong flow fraction:
\begin{equation}
  T_{7h} = T_{12h} - \epsilon_{\mathrm{IHX2}}\,
           \underbrace{x_{4h}}_{\text{\textbf{wrong}}}\,
           \mathrm{cpr}_{2}\left(T_{12h}-T_{13h}\right).
\end{equation}
$x_{4h}$ is the vapour \emph{quality} at the evaporator inlet. A quality is not
a flow split, and a flow split is what belongs there. The stream crossing IHX-2
on the vapour side is the suction stream --- what left the flash separator as
\emph{liquid}, passed through the evaporator, and is on its way to the low-stage
compressor --- which is $1-x_{8h}$. IHX-1, one screen earlier, correctly uses
$x_{8h}$, the flashed vapour.

The two numbers are close enough to look plausible and different enough to break
the first law:

\begin{center}
\begin{tabular}{@{}l r@{}}
\toprule
per \si{\kilogram} of high-stage flow & \si{\kilo\joule\per\kilogram} \\
\midrule
IHX-2, heat given by the liquid    & 9.19 \\
IHX-2, heat taken by the vapour    & 21.76 \\
\midrule
condenser $h_{2h}-h_{3h}$          & 271.43 \\
compressor work $+$ evaporator heat & 258.33 \\
\textbf{energy from nowhere}       & \textbf{13.11} \;(\SI{4.8}{\percent}) \\
\bottomrule
\end{tabular}
\end{center}

\medskip
\noindent\textbf{The repair.} The effectiveness loop is gone, replaced by the
enthalpy balance it was approximating:
\begin{align}
  h_{12h} &= h_{3h}  - x_{8h}\left(h_{11h}-h_{10h}\right)
    && \text{IHX-1, vapour} = x_{8h}, \notag\\
  h_{7h}  &= h_{12h} - \left(1-x_{8h}\right)\left(h_{1h}-h_{13h}\right)
    && \text{IHX-2, vapour} = 1-x_{8h}, \label{eq:ihx-balance}\\
  x_{8h}  &= \frac{h_{7h}-h_{9h}}{h_{10h}-h_{9h}}
    && \text{separator, } h_{8h}=h_{7h}. \notag
\end{align}
$x_{8h}$ appears on both sides, so Eq.~\eqref{eq:ihx-balance} is a fixed point.
It is linear and strongly contracting --- the gain is $0.013$ --- and reaches
$10^{-12}$ in a handful of iterations. State $1h$ is fixed independently by the
isentropic compression back from $2h$, so $\epsilon_{\mathrm{IHX2}}$ becomes an
\textbf{output}: the effectiveness the regenerator would have to have. It is
reported, and a value above unity would mean the component cannot exist. At the
design point it is $0.2338$.

\begin{center}
\begin{tabular}{@{}l r r@{}}
\toprule
 & v0.8 & v0.9 \\
\midrule
$x_{8h}$                 & 0.4160 & 0.3778 \\
$\mathrm{COP}$           & 2.9563 & 2.8868 \\
cycle residual [\si{\kilo\joule\per\kilogram}]
                         & 13.11 & \num{1.7e-12} \\
\bottomrule
\end{tabular}
\end{center}

\medskip
\noindent\textbf{Why it survived.} Nothing ever summed the cycle. The COP
expression is written in terms of $x_{8h}$ and is correct as written; it simply
consumed an $x_{8h}$ produced by an inconsistent network. An energy balance over
the whole cycle is one line, and it would have caught this on the first day. It
is now an acceptance test, and so is the IHX-2 balance on its own.

\medskip
\noindent
Both defects in this subsection and the last were found by a reader following
the cycle code by hand --- not by any of the internal consistency machinery of
\S\ref{sec:defects:verification}, which was looking only at the store.

%---------------------------------------------------------------------
\subsection{The heat-pump evaporator had no approach at all}
\label{sec:defects:hpe}

The two previous subsections concern quantities that were computed wrongly.
This one concerns a quantity that was never computed.

Two independent fields set the two ends of the same exchanger:
\begin{equation}
  T_{4a} = T_{\mathrm{source}} - \Delta T_{3A,4A},
  \qquad
  T_{13h} = T_{\mathrm{source}} - \Delta T_{3A,13H},
  \label{eq:hpe-old}
\end{equation}
so the approach between the source outlet and the evaporating temperature was
their \emph{difference}, and nothing anywhere checked its sign:

\begin{center}
\begin{tabular}{@{}r r r r r l@{}}
\toprule
$\Delta T_{3A,4A}$ & $\Delta T_{3A,13H}$ & source out & evaporating & approach
 & reported $\mathrm{COP}$ \\
\midrule
10 & 10 & \SI{50.0}{\celsius} & \SI{50.0}{\celsius} & $+0.0$ & 2.8868 \\
 8 & 10 & \SI{52.0}{\celsius} & \SI{50.0}{\celsius} & $+2.0$ & 2.8868 \\
12 & 10 & \SI{48.0}{\celsius} & \SI{50.0}{\celsius} & $\mathbf{-2.0}$ & 2.8868 \\
15 & 10 & \SI{45.0}{\celsius} & \SI{50.0}{\celsius} & $\mathbf{-5.0}$ & 2.8868 \\
\bottomrule
\end{tabular}
\end{center}

\noindent
The defaults put the approach at exactly \emph{zero} --- infinite area. Rows
three and four are second-law violations, and the model reports the same
$\mathrm{COP}$ for all four.

\medskip
\noindent\textbf{Why no test could have caught it.} $T_{4a}$ was written into
the state dictionary at one line and \textbf{never read again}. The source
stream's energy balance was never closed, so there was no exchanger there to
pinch: the evaporating temperature was declared, the source outlet was declared,
and the two were never required to be consistent. A crossing could not produce a
symptom because nothing downstream depended on the quantity that crossed.

That generalises, and it is the reason this subsection exists:

\begin{center}
\emph{A parameter that influences no output cannot be validated by any amount of
checking the outputs.}
\end{center}

\noindent
$\Delta T_{3A,4A}$ was inert. Energy closure, the CSS identity
(Eq.~\eqref{eq:css-unity}), the frozen fixture and the equivalence tests of
\S\ref{sec:defects:verification} all passed with it set to a value that violates
the second law. The guard added at v0.8 was worse than useless: it warned when
the two fields were \emph{equal} --- the benign case --- and was silent when
they were reversed.

\medskip
\noindent\textbf{The repair.} The evaporating temperature is now derived rather
than declared,
\begin{equation}
  \boxed{\;T_{13h} \;=\; T_{\mathrm{source}} - \Delta T_{3A,4A}
          - \Delta T_{\mathrm{pinch,HPE}}\;}
  \qquad \Delta T_{\mathrm{pinch,HPE}} = \SI{5}{\kelvin},
  \label{eq:hpe-new}
\end{equation}
so the approach is positive by construction. $\Delta T_{3A,13H}$ is retired, and
\code{validate\_case} raises if a Case still names it rather than ignoring it
silently. $\Delta T_{3A,4A}$ is no longer inert: it now moves the evaporating
temperature, and with it the $\mathrm{COP}$.

\medskip
\noindent\textbf{Why an end approach is sufficient here, when
\S\ref{sec:defects:orcpinch} needed a composite-curve search.} The cold stream
in this evaporator is isothermal with \emph{no preheat section}: state $4h$
enters two-phase and $13h$ leaves saturated, so its composite is flat and has no
kink. The hot stream is monotonic. The gap is therefore monotonic in $Q$ and its
minimum lies necessarily at the cold end, which is exactly what
$\Delta T_{\mathrm{pinch,HPE}}$ names. The ORC evaporator needs the search
precisely because its liquid preheat puts a kink \emph{before} the plateau, and
that is what moves the pinch into the interior.

\medskip
\noindent\textbf{What it costs.}

\begin{center}
\begin{tabular}{@{}l r r@{}}
\toprule
 & before & after \\
\midrule
evaporating temperature & \SI{50.0}{\celsius} & \SI{45.0}{\celsius} \\
$\mathrm{COP}$          & 2.8868 & 2.7594 \\
$\eta_{\mathrm{RTE}}$   & 0.3140 & 0.3003 \\
\bottomrule
\end{tabular}
\end{center}

\noindent
Everything on the discharge side is untouched to the last digit --- $\eta_{T}$,
the net electric energy per well, $\Nwells$, the CSS deviation, the thermal
energy per well and the storage density --- because the heat pump appears only
in the charging branch of the budget.

\medskip
\noindent\textbf{Still open.} The source stream is specified but not sized:
nothing tracks its mass flow, so the cost of cooling it by \SI{10}{\kelvin}
rather than \SI{5}{\kelvin} remains invisible. Closing that would make
$\Delta T_{3A,4A}$ a genuine trade rather than a free choice.

%---------------------------------------------------------------------
\subsection{Version history}

\begin{longtable}{@{}l l p{0.60\linewidth}@{}}
\toprule
version & date & change \\
\midrule
\endhead
\textbf{0.9} & 2026-09 &
  Both thermodynamic cycles repaired. The ORC evaporating temperature is set by
  a composite-curve pinch on the whole evaporator rather than by its hot end
  (\S\ref{sec:defects:orcpinch}); the crossing had been \SI{-29.8}{\kelvin}.
  The heat-pump IHX-2 balance is closed on enthalpy with the correct flow split
  $1-x_{8h}$ in place of the evaporator-inlet quality $x_{4h}$
  (\S\ref{sec:defects:ihx2}); the cycle had been out of balance by
  \SI{4.8}{\percent}. The heat-pump evaporator is closed on an approach that
  exists: the evaporating temperature is derived from the source outlet rather
  than declared beside it (\S\ref{sec:defects:hpe}), where the two free fields
  had made the approach exactly zero and could have made it negative in
  silence. $\eta_{\mathrm{ORC}}$ $0.2320\to0.1738$, $\mathrm{COP}$
  $2.9563\to2.7594$, $\Nwells$ $11.632\to15.659$, $\eta_{\mathrm{RTE}}$
  $0.4367\to0.3003$. \emph{Every} storage quantity is unchanged to the last
  digit carried --- CSS deviation, thermal energy per well, cycled fraction,
  storage density --- which is the evidence that the two corrections are
  confined to the cycles. New guard rails in \code{validate\_case}: boiling of
  the charging water, the zero approaches at the heat-pump source outlet and in
  the ORC regenerator, and a report whenever the pinch forces the boiling
  temperature down; \code{DT\_3A\_13H} is retired and raises if named. \\
0.5--0.8 & 2026-09 &
  Cyclic steady state replaces first-cycle sizing (\S\ref{sec:cyclic:reform});
  two recording defects in the history corrected; the discharge closed on its
  inlet rather than its outlet (\S\ref{sec:cyclic:closure}); the conduction
  shell taken on the correct side of the front
  (\S\ref{sec:network:shell}); the single-phase branches given a proper
  mean-to-surface bulk resistance rather than zero or the full annulus. \\
\textbf{0.4a} & 2026-09 &
  No physics. \S\ref{sec:front:enthalpy} rewritten as a teaching treatment of
  Formulation~C. Reporting rule added (\S\ref{sec:hydraulics:capacity}):
  $N_{\mathrm{capacity}}$ is a lower bound and must not be read across a sweep
  alone. Melt-front limit corrected from the borehole wall to the cell radius,
  Eq.~\eqref{eq:rcell}, and H3 proximity now reported
  (\S\ref{sec:defects:h3}); the fin count and the charging window follow from
  it. The segment equation is derived in full (\S\ref{sec:front:segeq}).
  Hypothesis \textbf{H9} (perfect azimuthal insulation between cascade layers)
  stated explicitly and quantified (\S\ref{sec:defects:h9}). Sizing vocabulary
  reduced to the two independent numbers it always had
  (\S\ref{sec:hydraulics:twoq}): \code{N\_inventory} and \code{N\_heat} are
  retired and now raise --- the first because it was a silent alias for
  $N_{\mathrm{capacity}}$ that had meant a \emph{different} quantity in v0.3,
  the second because it was one of three labels for the rate criterion. \\
0.4 & 2026-09 &
  Melt front carries \emph{enthalpy} rather than melted area
  (\S\ref{sec:front:enthalpy}): superheat above $\Tm$, subcool below it, no
  clipping, local melt fraction confined to $[0,1]$ by construction. Branchwise
  exact time integration replaces explicit Euler, which is unstable once the PCM
  has finite heat capacity. Sizing moves to a capacity criterion
  (\S\ref{sec:hydraulics:capacity}). $\Nwells$ $12.74\to11.57$; spread in local
  melt fraction $1.321\to0.078$. The v0.3 lower-bound comparison is corrected
  (\S\ref{sec:hydraulics:ideal}); it compared inconsistent bases. \\
0.3 & 2026-09 &
  Wall conductivity corrected to steel during charging
  (\S\ref{sec:defects:kw}); fin metal excluded from the melted PCM area
  (\S\ref{sec:front:finmetal}). \\
0.2 & 2026-08 &
  Melt front reformulated by energy balance (\S\ref{sec:front:marched}); melt
  state carried across the cycle; two-constraint sizing; latent inventory made
  density-consistent. \\
0.1 & --- &
  IHTC paper. Closed-form Stefan front, single sizing criterion, wall
  conductivity error. Reproduced exactly by the regression of
  \S\ref{sec:defects:verification}. \\
\bottomrule
\end{longtable}
